\chapter{SABR en taux et cube de volatilité}\label{ch:rc:sabr-in-rates-and-the-volatility-cube}

En 2015, un desk de taux euro demande à son modèle de smile la volatilité
d'une swaption receveuse de strike moins vingt points de base. Le modèle est SABR, le
standard du marché depuis 2002, dans la forme lognormale qu'utilisent tous les desks~: sa
formule prend $\ln(F/K)$ et $(FK)^{(1-\beta)/2}$, et avec un strike négatif elle
n'a rien à renvoyer. En quelques mois, tout le marché euro s'est rallié à l'une de
deux corrections~: ajouter une constante à chaque taux avant d'appliquer la formule,
ou utiliser la version du modèle où les taux sont normaux. Chacune ajuste les
strikes cotés~; elles divergent sur les strikes que personne ne cote. Ce chapitre
applique SABR (le livre~5, chapitre~11) aux taux, construit à partir de lui le cube
de volatilité des swaptions, nappe de smile après nappe de smile, l'interpole,
et regarde où il se brise~: aux strikes négatifs, et dans l'aile basse des échéances
longues, là où la densité de la formule devient négative.

\section{SABR en taux~: SABR normal et backbone}

Un taux swap à terme $F$, sous la probabilité annuité, suit la dynamique SABR
\[
dF_t = \alpha_t F_t^{\beta}\,dW_t,\qquad d\alpha_t = \nu\alpha_t\,dZ_t,\qquad
d\langle W,Z\rangle_t = \rho\,dt ,
\]
avec les paramètres $(\alpha,\beta,\rho,\nu)$ du livre~5. Les desks de taux fixent
$\beta$ par convention plutôt que par ajustement, car $\beta$ et $\rho$ inclinent
tous deux le smile et une seule tranche ne permet pas de les séparer~; le choix
fixe le backbone, c'est-à-dire la manière dont la volatilité à la monnaie évolue quand le forward bouge.

\begin{definition}[SABR normal]\label{def:rc:sabr-in-rates-and-the-volatility-cube:normal}
Le \emph{SABR normal}\index{SABR normal} est le SABR avec $\beta=0$~: $dF_t=\alpha_t\,dW_t$,
un taux normal compte tenu de sa volatilité et pouvant prendre n'importe quel signe. Sa
volatilité normale à la monnaie ne dépend pas du niveau du forward~:
le backbone est plat en termes normaux.
\end{definition}

\begin{proposition}[Le développement en volatilité normale du SABR normal]\label{prop:rc:sabr-in-rates-and-the-volatility-cube:normalexp}
Au premier ordre en l'échéance $T$ (Hagan et al., 2002), la volatilité implicite
normale du \omterm{def:rc:sabr-in-rates-and-the-volatility-cube:normal}{SABR normal} est
\[
\sigma_N(K) = \alpha\,\frac{\zeta}{x(\zeta)}\Bigl(1+\frac{2-3\rho^2}{24}\,\nu^2T\Bigr),
\quad \zeta = \frac{\nu}{\alpha}(F-K),\quad
x(\zeta) = \ln\frac{\sqrt{1-2\rho\zeta+\zeta^2}+\zeta-\rho}{1-\rho},
\]
avec $\zeta/x(\zeta)=1$ à la monnaie. Le smile ne dépend que de $F-K$~:
il se déplace avec le forward, strike par strike, en points de base.
\end{proposition}
\admitted\ (Hagan et al., 2002, le démontrent par perturbation singulière~; le livre~5,
chapitre~11, traite le cas lognormal.)

\omcode{code/firm/sabrcube/firm_sabrcube.py}{41}{61}{Développement lognormal de Hagan sur des taux décalés, et volatilité normale du SABR normal.}\label{lst:rc:sabr-in-rates-and-the-volatility-cube:hagan}

\section{Taux négatifs et modèles décalés}

\begin{definition}[Modèle lognormal décalé, SABR décalé]\label{def:rc:sabr-in-rates-and-the-volatility-cube:shifted}
Un \emph{modèle lognormal décalé}\index{modèle lognormal décalé} de décalage
$\zeta>0$ rend $F+\zeta$ lognormal, $dF_t=\sigma(F_t+\zeta)\,dW_t$, de sorte que le
taux est borné inférieurement par $-\zeta$~; les options se valorisent avec la formule de Black sur
$F+\zeta$ et $K+\zeta$. Le \emph{SABR décalé}\index{SABR décalé} applique la dynamique du
SABR à $F+\zeta$, $dF_t=\alpha_t(F_t+\zeta)^\beta\,dW_t$, et le développement lognormal
de Hagan au forward et au strike décalés.
\end{definition}

Le décalage est un plancher sur le taux modélisé, fixé à l'avance. Il doit se situer
sous tout taux que le marché peut raisonnablement atteindre.

\begin{definition}[Plancher effectif des taux]\label{def:rc:sabr-in-rates-and-the-volatility-cube:elb}
Le \emph{plancher effectif des taux}\index{plancher effectif des taux} d'un taux directeur est
le niveau en deçà duquel une banque centrale ne devrait pas baisser, parce que de nouvelles
baisses feraient sortir les liquidités du système bancaire, ou pénaliseraient davantage
les banques qu'elles n'assoupliraient les conditions. Les modèles de taux s'en servent pour borner
la distribution des taux, et le décalage d'un modèle décalé doit en dépasser la valeur absolue.
\end{definition}

\begin{dated}{2026-09}{Où le plancher a été franchi}\label{dat:rc:sabr-in-rates-and-the-volatility-cube:elb}
Le taux de dépôt de la BCE a été négatif de juin 2014 à juillet 2022, avec un point bas
de $-0.50~\%$ à partir de septembre 2019~; le taux directeur de la Banque nationale suisse
a été de $-0.75~\%$ à partir de janvier 2015 et est repassé au-dessus de zéro en septembre 2022~; celui de la Banque du Japon
a été de $-0.1~\%$ de 2016 jusqu'au 19 mars 2024.
\end{dated}

\begin{example}[Un smile, trois modèles]\label{ex:rc:sabr-in-rates-and-the-volatility-cube:neg}
Une option d'un an sur un swap euro de deux ans, de forward $-0.30~\%$, est cotée en
volatilités normales de 29.3, 26.5, 25.7, 27.8, 31.4 et 39.8 points de base aux
strikes 50 et 25 points de base sous le forward, à la monnaie, puis 25, 50 et
100 au-dessus (un smile engendré par le \omterm{def:rc:sabr-in-rates-and-the-volatility-cube:normal}{SABR normal}, à titre illustratif). Le \omterm{def:rc:sabr-in-rates-and-the-volatility-cube:shifted}{SABR décalé} avec $\beta=0.5$
l'ajuste à 0.77 points de base d'erreur quadratique moyenne avec un décalage de
1~\%, à 0.11 avec un décalage de 3~\%~; le \omterm{def:rc:sabr-in-rates-and-the-volatility-cube:normal}{SABR normal} l'ajuste exactement. Interrogé
sur une receveuse de strike $-0.90~\%$, soixante points de base sous le forward et hors des
cotations, il répond 0.177, 0.280 et 0.295 points de base d'annuité
(\cref{fig:rc:sabr-in-rates-and-the-volatility-cube:neg})~: avec un décalage de 1~\%, le
taux modélisé ne peut pas descendre sous $-1~\%$, et la receveuse profonde coûte 37~\% de
moins qu'avec 3~\%.
\end{example}

\begin{omfigure}
\begin{tikzpicture}
\begin{axis}[width=12cm,height=5.8cm,
  xlabel={strike (\%)},ylabel={volatilité normale (pb)},
  tick label style={font=\small},label style={font=\small},
  xmin=-1,xmax=0.8,ymin=24,ymax=45,grid=major,grid style={gray!25},unbounded coords=jump,
  legend columns=4,legend style={font=\footnotesize,at={(0.5,-0.25)},anchor=north,draw=none,fill=none,/tikz/every even column/.append style={column sep=6pt}},
  table/col sep=comma]
\addplot[omThm,thick] table[x=k,y=s1]{figdata/rates-credit-risk/05-sabr-in-rates-and-the-volatility-cube/negsmile.csv};
\addplot[omDef,thick,dashed] table[x=k,y=s3]{figdata/rates-credit-risk/05-sabr-in-rates-and-the-volatility-cube/negsmile.csv};
\addplot[omProp,thick,dotted] table[x=k,y=nrm]{figdata/rates-credit-risk/05-sabr-in-rates-and-the-volatility-cube/negsmile.csv};
\addplot[only marks,mark=*,mark size=2pt,black] table[x=k,y=vol]{figdata/rates-credit-risk/05-sabr-in-rates-and-the-volatility-cube/negquotes.csv};
\legend{SABR décalé 1~\%,SABR décalé 3~\%,SABR normal,cotations}
\draw[gray,dashed] (axis cs:-0.9,24) -- (axis cs:-0.9,45) node[pos=0.93,right,font=\footnotesize,text=black] {$-0.90~\%$};
\end{axis}
\end{tikzpicture}

\omcaption{Un smile de taux négatifs (un an dans deux ans, taux forward $-0.30~\%$)
ajusté par trois modèles, en volatilité normale. Le décalage de 3~\% et le \omterm{def:rc:sabr-in-rates-and-the-volatility-cube:normal}{SABR normal}
passent par les cotations~; le décalage de 1~\% s'en écarte de jusqu'à 1.2~points de base et, en dessous, son
volatilité descend vers son plancher à $-1\%$. Données~: cotations synthétiques issues du \omterm{def:rc:sabr-in-rates-and-the-volatility-cube:normal}{SABR normal}~; le
tutoriel du chapitre.}\label{fig:rc:sabr-in-rates-and-the-volatility-cube:neg}
\end{omfigure}

\section{Construire le cube, smile par smile}

\begin{definition}[Matrice des swaptions]\label{def:rc:sabr-in-rates-and-the-volatility-cube:matrix}
La \emph{matrice des swaptions}\index{matrice des swaptions} d'une devise est le tableau des
volatilités des swaptions à la monnaie par échéance de l'option et par durée du swap sous-jacent~:
c'est la face à la monnaie du cube de volatilité (livre~2, chapitre~13).
\end{definition}

\begin{method}[Calibrer le cube]\label{met:rc:sabr-in-rates-and-the-volatility-cube:cube}
Pour chaque couple (échéance, durée) coté~: calculer le taux de swap forward et son annuité
sur la courbe multi-courbes (chapitre~2)~; convertir les cotations (volatilités normales,
ou prix) en la volatilité propre au modèle (volatilités de Black décalées pour le
\omterm{def:rc:sabr-in-rates-and-the-volatility-cube:shifted}{SABR décalé})~; fixer $\beta$ et le décalage pour l'ensemble de la devise~; ajuster $(\alpha,\rho,\nu)$ par
moindres carrés, avec $\alpha>0$ et $|\rho|<1$ imposés par reparamétrage~;
consigner l'erreur d'ajustement. Vérifier ensuite la régularité des paramètres le long de la grille~:
une section dont $\rho$ ou $\nu$ fait un saut est en général une mauvaise cotation.
\end{method}

\begin{omfigure}
\begin{tikzpicture}[x={(1cm,0cm)},y={(0.5cm,0.4cm)},z={(0cm,1cm)},font=\small]
\draw[omDef,thick] (0,0,0) -- (5,0,0) -- (5,4,0) -- (0,4,0) -- cycle;
\draw[omDef,thick] (0,0,3) -- (5,0,3) -- (5,4,3) -- (0,4,3) -- cycle;
\foreach \x/\y in {0/0,5/0,5/4,0/4} {\draw[omDef,thick] (\x,\y,0) -- (\x,\y,3);}
\fill[omThm!18] (0,0,1.5) -- (5,0,1.5) -- (5,4,1.5) -- (0,4,1.5) -- cycle;
\draw[omThm,thick] (0,0,1.5) -- (5,0,1.5) -- (5,4,1.5) -- (0,4,1.5) -- cycle;
\draw[omProp,very thick] (2.5,2,0) -- (2.5,2,3);
\node[anchor=north] at (2.5,0,0) {échéance de l'option};
\node[anchor=west] at (5.1,2.2,0) {durée du swap};
\node[anchor=east] at (0,0,1.5) {strike};
\node[text=omThm,anchor=west] (atm) at (5.9,4,1.1) {face à la monnaie~: la matrice des swaptions};
\node[text=omProp,anchor=west] (sec) at (5.9,4,3.0) {une section du smile~: un ajustement SABR};
\draw[omThm,thin] (atm.west) -- (4.2,3.2,1.5);
\draw[omProp,thin] (sec.west) -- (2.5,2,2.6);
\end{tikzpicture}

\omcaption{Le cube de volatilité~: un smile pour chaque couple (échéance, durée). La
face à la monnaie est la \omterm{def:rc:sabr-in-rates-and-the-volatility-cube:matrix}{matrice des swaptions}~; chaque ligne verticale est une section
calibrée avec ses propres paramètres SABR, $\beta$ et le décalage étant fixés pour
l'ensemble de la devise.}\label{fig:rc:sabr-in-rates-and-the-volatility-cube:cube}
\end{omfigure}

\begin{example}[Un cube synthétique en euros]\label{ex:rc:sabr-in-rates-and-the-volatility-cube:cube}
Seize sections (échéances de un, deux, cinq et dix ans~; durées de deux, cinq, dix
et trente ans), sept strikes allant de $-200$ à $+200$ points de base autour
du forward sur les courbes du chapitre~2, engendrées par le \omterm{def:rc:sabr-in-rates-and-the-volatility-cube:normal}{SABR normal} avec des paramètres
qui varient régulièrement avec l'échéance et la durée, plus un bruit à graine fixée de 0.3 point de base,
sont calibrées avec le \omterm{def:rc:sabr-in-rates-and-the-volatility-cube:shifted}{SABR décalé} ($\beta=0.5$, décalage 3~\%). Chaque section s'ajuste
avec un écart compris entre 0.21 et 0.78 point de base. La face à la monnaie
(\cref{fig:rc:sabr-in-rates-and-the-volatility-cube:matrix}) croît avec l'échéance et
décroît avec la durée~: 69.9~points de base pour un an dans deux ans, 74.0 pour dix dans
trente.
\end{example}

\begin{omfigure}
\begin{tikzpicture}
\begin{axis}[width=11cm,height=5.4cm,
  xlabel={échéance de l'option (années)},ylabel={volatilité normale ATM (pb)},
  tick label style={font=\small},label style={font=\small},
  xmin=0,xmax=10.5,ymin=50,ymax=95,grid=major,grid style={gray!25},
  legend columns=4,legend style={font=\footnotesize,at={(0.5,-0.25)},anchor=north,draw=none,fill=none,/tikz/every even column/.append style={column sep=6pt}},
  table/col sep=comma]
\addplot[omDef,thick,mark=*] table[x=e,y=n2]{figdata/rates-credit-risk/05-sabr-in-rates-and-the-volatility-cube/matrix.csv};
\addplot[omThm,thick,mark=square*] table[x=e,y=n5]{figdata/rates-credit-risk/05-sabr-in-rates-and-the-volatility-cube/matrix.csv};
\addplot[omProp,thick,mark=triangle*] table[x=e,y=n10]{figdata/rates-credit-risk/05-sabr-in-rates-and-the-volatility-cube/matrix.csv};
\addplot[omExo,thick,mark=diamond*] table[x=e,y=n30]{figdata/rates-credit-risk/05-sabr-in-rates-and-the-volatility-cube/matrix.csv};
\legend{durée 2 ans,durée 5 ans,durée 10 ans,durée 30 ans}
\end{axis}
\end{tikzpicture}

\omcaption{La \omterm{def:rc:sabr-in-rates-and-the-volatility-cube:matrix}{matrice des swaptions} du cube synthétique~: volatilité normale à la monnaie des
sections calibrées par échéance, une courbe par durée. Données~: le
tutoriel du chapitre.}\label{fig:rc:sabr-in-rates-and-the-volatility-cube:matrix}
\end{omfigure}

\section{Interpoler en échéance, en durée et en strike}

Le cube est coté sur une grille clairsemée et doit être disponible partout~: un swap de sept ans
démarrant dans trois ans, une swaption à un strike quelconque. La dimension des strikes est
traitée par la formule SABR elle-même~; l'échéance et la durée le sont en
interpolant les paramètres, et non les volatilités, de sorte que chaque smile interpolé
soit lui-même un smile SABR.

\begin{example}[Une section que personne ne cote]\label{ex:rc:sabr-in-rates-and-the-volatility-cube:missing}
Interpoler $(\alpha,\rho,\nu)$ bilinéairement en échéance et en durée entre les
échéances de deux et cinq ans et les durées de cinq et dix ans donne le smile trois ans dans
sept ans (forward 2.905~\%) à moins de 0.93 point de base du smile que les
paramètres générateurs impliquent, sur les sept strikes.
\end{example}

\begin{remark}[Ce que l'interpolation doit respecter]\label{rem:rc:sabr-in-rates-and-the-volatility-cube:interp}
L'interpolation linéaire de $\alpha$ en échéance est une règle grossière~: c'est la variance totale,
et non la volatilité, qui doit croître avec l'échéance, et $\alpha$ n'en est que le terme dominant.
Les desks interpolent en variance le long de l'échéance, gardent $\rho$ et $\nu$ réguliers, et
vérifient les sections interpolées pour l'arbitrage calendaire et l'arbitrage butterfly (livre~5, chapitre~7).
Les forwards proviennent de la courbe, jamais de l'interpolation.
\end{remark}

\section{Arbitrage dans les ailes}

La formule de Hagan est un développement, précis pour un $\nu^2T$ modéré et des strikes
pas trop éloignés du forward. Sa densité implicite, dérivée seconde par rapport au strike du
prix du call (Breeden--Litzenberger, livre~5, chapitre~7), peut devenir
négative dans l'aile des bas strikes des échéances longues~: un butterfly de
prix négatif.

\omcode{code/firm/sabrcube/firm_sabrcube.py}{136}{140}{Densité implicite d'une section, par différences secondes des prix.}\label{lst:rc:sabr-in-rates-and-the-volatility-cube:density}

\begin{example}[Dix ans sur dix ans]\label{ex:rc:sabr-in-rates-and-the-volatility-cube:density}
Une option à dix ans sur un swap à dix ans à 2.50~\%, cotée de $-200$ à $+200$
points de base autour du forward (86 à 96 points de base de volatilité normale),
est calée par un \omterm{def:rc:sabr-in-rates-and-the-volatility-cube:shifted}{SABR décalé} (shifted SABR), avec un décalage de 1~\% (erreur de 1.28 points de base) et de
3~\% (0.35). Avec le décalage de 1~\%, la densité est négative de $-0.95~\%$ à
$+0.35~\%$, atteignant $-158$~; avec 3~\%, elle reste positive jusqu'à $-2.15~\%$
(la \cref{fig:rc:sabr-in-rates-and-the-volatility-cube:density}). Un portefeuille de
swaptions receveuses à bas strikes valorisées avec l'ajustement à 1~\% portent un arbitrage contre les
butterflies du marché.
\end{example}

\begin{omfigure}
\begin{tikzpicture}
\begin{axis}[width=12cm,height=5.6cm,
  xlabel={taux à l'échéance (\%)},ylabel={densité},
  tick label style={font=\small},label style={font=\small},
  xmin=-1,xmax=7,ymin=-165,ymax=30,grid=major,grid style={gray!25},unbounded coords=jump,
  legend columns=3,legend style={font=\footnotesize,at={(0.5,-0.25)},anchor=north,draw=none,fill=none,/tikz/every even column/.append style={column sep=8pt}},
  table/col sep=comma]
\addplot[omThm,thick] table[x=k,y=s1]{figdata/rates-credit-risk/05-sabr-in-rates-and-the-volatility-cube/density.csv};
\addplot[omDef,thick,dashed] table[x=k,y=s3]{figdata/rates-credit-risk/05-sabr-in-rates-and-the-volatility-cube/density.csv};
\addplot[omProp,thick,dotted] table[x=k,y=nrm]{figdata/rates-credit-risk/05-sabr-in-rates-and-the-volatility-cube/density.csv};
\legend{SABR décalé 1~\%,SABR décalé 3~\%,SABR normal (modèle générateur)}
\draw[gray] (axis cs:-1,0) -- (axis cs:7,0);
\end{axis}
\end{tikzpicture}

\omcaption{Densités implicites du taux du swap à dix ans sur dix ans issues de
trois ajustements sur les mêmes cotations. La densité du décalage à 1~\% devient fortement négative
juste au-dessus de son plancher~; celle du décalage à 3~\% reste positive sur l'intervalle représenté.
Données~: cotations synthétiques issues du \omterm{def:rc:sabr-in-rates-and-the-volatility-cube:normal}{SABR normal}~; le
tutoriel du chapitre.}\label{fig:rc:sabr-in-rates-and-the-volatility-cube:density}
\end{omfigure}

\begin{remark}[Corrections pour les ailes]\label{rem:rc:sabr-in-rates-and-the-volatility-cube:free}
Trois sont retenues~: un décalage plus grand (qui place le problème sous tout strike pertinent)~;
la résolution exacte ou numérique de la dynamique SABR, au lieu de la développer (la
densité est alors positive par construction)~; et le SABR à frontière libre de
Antonov, Konikov et Spector (2015), qui remplace $F^\beta$ par $|F|^\beta$ et
laisse le taux franchir zéro sans décalage, avec un prix exact pour une
corrélation nulle. Quel que soit le modèle, un contrôle de la densité sur les strikes
qu'un portefeuille contient est un test de validation standard (du \cref{ch:rc:model-risk-and-validation}).
\end{remark}

\section{Tutoriel~: un cube de volatilité à partir de cotations}\label{tut:rc:sabr-in-rates-and-the-volatility-cube:tutorial}

\begin{tutorial}
\textbf{Objectif.} Caler un smile de taux négatifs avec trois modèles, calibrer et
interpoler un cube, et vérifier sa densité. \textbf{État final~:}
\cref{fig:rc:sabr-in-rates-and-the-volatility-cube:neg,fig:rc:sabr-in-rates-and-the-volatility-cube:matrix,fig:rc:sabr-in-rates-and-the-volatility-cube:density}.
\begin{tutsteps}
\item \textbf{Un smile}~: \texttt{neg\_fits()} calibre le \omterm{def:rc:sabr-in-rates-and-the-volatility-cube:shifted}{SABR décalé} à 1~\% et
3~\% et le \omterm{def:rc:sabr-in-rates-and-the-volatility-cube:normal}{SABR normal}~; \texttt{deep\_receiver()} valorise la swaption receveuse à $-0.90\%$.
\item \textbf{Le cube}~: \texttt{build\_cube()} renvoie les sections calibrées et
leurs erreurs~; \texttt{matrix\_table()} la face à la monnaie.
\item \textbf{Interpoler}~: \texttt{missing\_section(3, 7)}.
\item \textbf{Densité}~: \texttt{density\_table()}~; \texttt{fig\_rc\_sabrcube.py}
écrit les graphiques.
\end{tutsteps}
\textbf{Pour aller plus loin.} Caler avec $\beta=0$ et $\beta=1$ pour un décalage de 3~\% et
comparer $\rho$~; interpoler la variance totale plutôt que $\alpha$ le long des échéances et
mesurer l'erreur sur la section manquante.
\end{tutorial}

\section{Construction~: le cube de swaptions}\label{bld:rc:sabr-in-rates-and-the-volatility-cube:sabrcube}

\begin{build}
\bfield{Rôle} La surface de volatilité de toutes les options de taux du portefeuille~:
les swaptions, quelle que soit l'échéance, le tenor et le strike, la réplication CMS (chapitre~6),
les cibles de calibration bermudiennes (chapitres 7 à 9).
\bfield{Interface} \texttt{Sabr(alpha, beta, rho, nu, shift)}~;
\texttt{hagan\_lognormal}, \texttt{hagan\_normal\_beta0}, \texttt{price},
\texttt{normal\_vol}~; \texttt{calibrate\_section(f, t, strikes, normal\_vols,
beta, shift, model)}~; \texttt{density}~; \texttt{Cube(expiries, tenors, sections,
forwards)} avec \texttt{params}, \texttt{forward}, \texttt{normal\_vol}.
\bfield{Règles} Cotations en volatilité normale~; $\beta$ et le décalage fixés par
devise~; paramètres interpolés, forwards issus de la courbe~; numpy uniquement.
\bfield{Tests d'acceptation} \texttt{code/firm/sabrcube/tests/}~: la limite à la monnaie
du développement~; le \omterm{def:rc:sabr-in-rates-and-the-volatility-cube:normal}{SABR normal} sans vol-of-vol est Bachelier~; le
calibrateur retrouve les paramètres injectés~; payeur moins receveur vaut le forward~;
une densité bénigne est positive et s'intègre à un~; le cube interpole
les paramètres de façon bilinéaire.
\bfield{Pour aller plus loin} Interpolation en variance totale~; une EDP ou un SABR à intégration exacte
pour des densités positives~; calibration aux prix avec des pondérations cours acheteur--cours vendeur.
\end{build}

\begin{omsources}
\item P.\ Hagan, D.\ Kumar, A.\ Lesniewski et D.\ Woodward, «~Managing smile
risk~», \emph{Wilmott Magazine}, 2002.
\item A.\ Antonov, M.\ Konikov et M.\ Spector, «~The free boundary SABR:
natural extension to negative rates~», \emph{Risk}, 2015.
\item BCE, \emph{Key ECB interest rates}.
\item L.\ Andersen et V.\ Piterbarg, \emph{Interest Rate Modeling}, volume 2,
Atlantic Financial Press, 2010, chapitre 16.
\end{omsources}

\section{Exercices}

\begin{exercise}[$\star$]\label{exo:rc:sabr-in-rates-and-the-volatility-cube:1}
Un forward vaut $-0.30~\%$ et un strike $-0.50~\%$. Dire si le SABR lognormal standard permet de valoriser
l'option~? Déterminer le plus petit décalage qui le rend possible~?
\end{exercise}

\begin{exercise}[$\star$]\label{exo:rc:sabr-in-rates-and-the-volatility-cube:2}
Dans le \omterm{def:rc:sabr-in-rates-and-the-volatility-cube:normal}{SABR normal}, le forward monte de 20 points de base d'un jour à l'autre et les
paramètres restent inchangés. Déterminer ce qu'il advient de la volatilité normale à la monnaie,
et de la volatilité du strike qui était à la monnaie la veille~?
\end{exercise}

\begin{exercise}[$\star$]\label{exo:rc:sabr-in-rates-and-the-volatility-cube:3}
Pourquoi les desks de taux fixent-ils $\beta$ au lieu de le calibrer~?
\end{exercise}

\begin{exercise}[$\star\star$]\label{exo:rc:sabr-in-rates-and-the-volatility-cube:4}
Avec le \omterm{def:rc:sabr-in-rates-and-the-volatility-cube:normal}{SABR normal} pour $\alpha=25$ points de base, $\rho=0.1$, $\nu=0.6$ et $T=1$,
calculer la volatilité normale à la monnaie.
\end{exercise}

\begin{exercise}[$\star\star$]\label{exo:rc:sabr-in-rates-and-the-volatility-cube:5}
Dans l'\cref{ex:rc:sabr-in-rates-and-the-volatility-cube:neg}, pourquoi le décalage de 1~\%
s'ajuste-t-il moins bien aux strikes cotés et diffère-t-il aussi dans l'aile~?
\end{exercise}

\begin{exercise}[$\star\star$]\label{exo:rc:sabr-in-rates-and-the-volatility-cube:6}
Lire la \cref{fig:rc:sabr-in-rates-and-the-volatility-cube:density}. Décrire le
butterfly qui exploiterait l'ajustement à 1~\% près de 0~\%, et expliquer pourquoi ce n'est pas
de l'argent gratuit en pratique.
\end{exercise}

\begin{exercise}[$\star\star\star$]\label{exo:rc:sabr-in-rates-and-the-volatility-cube:7}
\emph{Programmation.} Parmi les décalages de 0.85~\%, 0.90~\%, 0.95~\% et 1~\%, trouver le plus petit
pour lequel le smile calibré pour un an dans deux ans admet une densité positive
de $-0.80~\%$ à $+1~\%$.
\end{exercise}

\begin{exercise}[$\star\star\star$]\label{exo:rc:sabr-in-rates-and-the-volatility-cube:8}
\emph{Trouver la faille.} «~Pour combler la section trois dans sept, nous interpolons les
volatilités normales cotées strike par strike entre les sections voisines.~»
Qu'est-ce qui ne va pas~?
\end{exercise}

\section{Problème~: Le strike à moins quatre-vingt-dix}

\begin{problem}[{Problème du week-end --- un client demande une swaption receveuse profonde}]\label{pb:rc:sabr-in-rates-and-the-volatility-cube:1}
Durant une année de taux négatifs, un fonds de pension demande à un desk une swaption receveuse~:
d'un an dans deux ans, sur 500~millions d'euros, avec un strike de $-0.90~\%$, pour se protéger
contre une baisse des taux bien plus forte. Le forward vaut $-0.30~\%$, l'annuité
vaut 2.0, et le marché cote les six strikes de
l'\cref{ex:rc:sabr-in-rates-and-the-volatility-cube:neg}.

\medskip\noindent\textbf{Partie I --- Les cotations.}
\begin{enumerate}
\item À quelle distance le strike est-il du forward, et de la cotation la plus basse~?
\item Pourquoi le SABR lognormal standard ne peut-il pas le valoriser~?
\item Quels ajustements sont possibles, et que suppose chacun sur le taux le plus bas~?
\item Donner les erreurs d'ajustement des trois modèles.
\item Lequel ajuste le mieux les strikes cotés~?
\end{enumerate}

\medskip\noindent\textbf{Partie II --- Les prix.}
\begin{enumerate}[resume]
\item Donner la prime en euros sous chaque modèle.
\item De combien les décalages de 1~\% et de 3~\% divergent-ils~?
\item Pourquoi le prix du décalage de 1~\% est-il le plus bas~?
\item À quel prix le desk préférerait-il vendre, et à quel prix le client préférerait-il
acheter~?
\item Quel est le plus petit décalage dont la densité est positive à partir de $-0.80~\%$~?
\end{enumerate}

\medskip\noindent\textbf{Partie III --- Risque.}
\begin{enumerate}[resume]
\item Le desk vend au prix à 3~\% et se couvre avec le
receveur coté à $-0.80~\%$. Quel risque subsiste-t-il~?
\item Comment la réponse change-t-elle si la BCE baisse à $-0.75~\%$~?
\item Pourquoi le décalage lui-même est-il un paramètre de modèle qui exige une réserve~?
\item Que suppose le \omterm{def:rc:sabr-in-rates-and-the-volatility-cube:normal}{SABR normal} sur le niveau auquel les taux peuvent descendre~?
\item Comment tester le modèle retenu avant de coter~?
\end{enumerate}

\medskip\noindent\textbf{Partie IV --- Jugement.}
\begin{enumerate}[resume]
\item Pourquoi le marché de l'euro s'est-il fixé sur des décalages de quelques pour cent plutôt que
le plus petit qui fonctionne~?
\item Un décalage unique pour l'ensemble du cube est-il raisonnable~?
\item Que dire au client sur l'intervalle de prix~?
\item Énoncer le \emph{résultat nommé}~: la prime sous les décalages de 1~\% et de 3~\%, et
le plus petit décalage admissible.
\item En une phrase~: qu'est-ce qui détermine le prix d'un strike que personne ne cote~?
\end{enumerate}
\end{problem}

\section{Questions d'entretien}

\begin{interviewq}[$\star$ \iqroles{trader, chercheur}]\label{iq:rc:sabr-in-rates-and-the-volatility-cube:1}
Comment le marché des taux a-t-il valorisé les options lorsque les taux sont devenus négatifs~?
\end{interviewq}

\begin{interviewq}[$\star\star$ \iqroles{chercheur}]\label{iq:rc:sabr-in-rates-and-the-volatility-cube:2}
Que contrôle $\beta$ dans le SABR, et pourquoi est-il en général fixé~?
\end{interviewq}

\begin{interviewq}[$\star\star$ \iqroles{chercheur, développeur}]\label{iq:rc:sabr-in-rates-and-the-volatility-cube:3}
Comment construiriez-vous un cube de swaptions à partir des cotations des courtiers~?
\end{interviewq}

\begin{interviewq}[$\star\star$ \iqroles{trader}]\label{iq:rc:sabr-in-rates-and-the-volatility-cube:4}
Le forward monte de dix points de base. Qu'advient-il du smile de volatilité sous le
\omterm{def:rc:sabr-in-rates-and-the-volatility-cube:normal}{SABR normal} et sous le SABR lognormal avec $\beta=1$~?
\end{interviewq}

\begin{interviewq}[$\star\star\star$ \iqroles{chercheur, risque}]\label{iq:rc:sabr-in-rates-and-the-volatility-cube:5}
La formule de Hagan donne une densité négative. Où, pourquoi, et que faites-vous~?
\end{interviewq}

\begin{interviewq}[$\star\star\star$ \iqroles{développeur}]\label{iq:rc:sabr-in-rates-and-the-volatility-cube:6}
Votre calibration du cube échoue chaque matin sur une section. Comment procédez-vous
pour enquêter~?
\end{interviewq}
